\section{Does the Depth Boundary Generalise Beyond Vision?}
\label{sec:crossmodal}

The boundary in Section~\ref{sec:transport_results} is a claim about where a
language model stops reading a non-text modality, not about images in
particular. We therefore applied the same masking protocol, unchanged,
to an audio--language model and a video--language model, and then asked
whether the two policies follow. Each modality is represented by a
single checkpoint, so this section reports transfer of a measurement,
not a multi-model claim; the results in Section~\ref{sec:results} remain
the basis for the paper's conclusions.

\subsection{The Boundary Is Modality-General; Its Depth Is Not}

Table~\ref{tab:crossmodal_boundary} reports the protocol applied to
three modalities. In each case masking every text row's attention to
the modality columns from layer $\ell$ onward leaves the output
unchanged once $\ell$ passes a fixed depth, and the all-layer control
moves the output substantially.

\begin{table}[tb]
\centering\small
\caption{The masking protocol of Section~\ref{sec:depth_boundary} applied
unchanged to three modalities. ``Completion'' is the shallowest suffix
mask whose mean output KL falls below $0.03$; ``peak'' is the depth of
the largest single-layer KL. Depths are fractions of the backbone's
layer count. Audio and video use the modality-used subset, matching the
convention used for vision. Each modality is one checkpoint.}
\label{tab:crossmodal_boundary}
\begin{tabular}{llrrr}
\toprule
Modality & Backbone & Layers & Completion & Peak \\
\midrule
Image & Qwen3-VL-2B   & 28 & $0.57$ & $0.39$ \\
Audio & Qwen2-Audio-7B & 32 & $0.75$ & $0.38$ \\
Video & Qwen3-VL-2B   & 28 & $0.57$ & $0.43$ \\
\bottomrule
\end{tabular}
\end{table}

Two observations follow. First, transport completes at a well-defined
depth in every modality tested: the phenomenon is a property of how
these backbones consume non-text tokens, not of images. Second, the
completion depth is not constant. Audio completes at $0.75$ rather
than $0.57$, which falsified the depth we had registered in advance for
that modality. The depth of peak single-layer sensitivity, by contrast,
is nearly invariant at $0.38$--$0.43$. The layer at which the model
reads a modality most strongly is stable; the layer by which it has
finished reading is not.

\subsection{What AVR Requires: an Information Knob, Not a Token Knob}

Section~\ref{sec:avr_results} established that AVR's gain equals the
accuracy difference between its high- and low-budget endpoints. Read as
a precondition rather than a diagnostic, this identity states when AVR
can pay at all: only when raising the budget restores information that
the low budget destroyed. Table~\ref{tab:crossmodal_headroom} applies
that test.

\begin{table}[tb]
\centering\small
\caption{Headroom, measured against the strongest equal-budget
reference in each setting, and the AVR gain it predicts. An asterisk
marks a 95\% interval excluding zero. Intervals are clustered by
source item. Video CIs are clustered by base video; audio uses the
first 180\,s of each clip.}
\label{tab:crossmodal_headroom}
\begin{tabular}{llrr}
\toprule
Setting & Budget knob & Headroom & AVR $-$ reference \\
\midrule
Audio, 10\,s clips  & time-stretch      & $-3.8$          & $-5.0$ \\
Audio, 180\,s clips & native encoding   & $+2.3^*$        & --- \\
Video, TempCompass  & frames $4\to32$   & $+4.9^*$        & $+3.4$ \\
\bottomrule
\end{tabular}
\end{table}

On short audio clips the only available knob is time-stretching, which
lengthens the token sequence without adding acoustic detail. Headroom is
negative and AVR does not pay, exactly as the identity predicts. On long
clips the encoder must compress, so encoding natively recovers
information the compressed input lost, and headroom becomes positive.
On video the frame count is a genuine resolution knob and headroom is
again positive.

Positive headroom is necessary but not sufficient. AVR spends its saving
on a larger budget, and the size of that saving is set by the boundary
itself. With $L$ layers, a cut at layer $P$, and a retained fraction
$k$, the affordable multiplier over the reference budget is
\begin{equation}
\kappa = \frac{L}{P + k\,(L-P)} .
\label{eq:kappa}
\end{equation}
A deeper boundary leaves fewer layers over which to save. For video,
$P/L = 0.57$ and $k = 0.1$ give $\kappa = 1.63$; for audio, $P/L=0.75$
gives $\kappa = 1.29$. The video headroom of $4.9$ points was measured
across an eightfold span, and at the affordable $1.5\times$ the measured
AVR gain is $+3.4$ points with an interval that includes zero --- which
matches the $+3.2$ points of headroom available at that same budget,
with pruning itself costing $+0.2$ points. The policy collects what is
available to it; the boundary determines how much that is.
Equation~\ref{eq:kappa} makes this quantitative rather than incidental:
the same measurement that locates the read-out also bounds the budget a
depth-aware policy can reinvest.

\subsection{What Placement Requires: Room Above the Reference}

For DWA the corresponding precondition is that some crop is better than
no crop. We test this with an oracle that selects among candidate
windows using the label, which upper-bounds any label-free placer. The
oracle must be compared against a reference that is selected the same
way: choosing the best of $K$ draws beats a single draw even when no
draw is better on average. On video, this selection effect alone is
worth $3.7$ points of the $9.1$ points an uncontrolled comparison
reports.

\begin{table}[tb]
\centering\small
\caption{Placement ceilings. The ceiling compares a label-selected crop
against a label-selected reference at the same budget, so it isolates
placement from selection. ``Incumbent'' is the block-mean read-out
against a random window.}
\label{tab:crossmodal_ceiling}
\begin{tabular}{lrr}
\toprule
Setting & Ceiling (bias-matched) & Incumbent $-$ random \\
\midrule
Audio, MMAU        & $-3.3$   & $+5.0^*$ \\
Video, TempCompass & $+5.5^*$ & $-0.7$ \\
\bottomrule
\end{tabular}
\end{table}

The two modalities fail the precondition in opposite ways
(Table~\ref{tab:crossmodal_ceiling}). In audio the label-free statistic
carries real signal, but the best window chosen with the label is still
worse than not cropping: placement has nothing to reach, so no placer
was fitted. In video there is genuine room, but the block-mean statistic
is indistinguishable from choosing at random. Refitting the read-out of
Section~\ref{sec:dwa_method} on video, with folds held out by source video,
reproduces the in-sample fit but does not generalise, and remains below
a budget-matched reference. What fails in video is not the value of
placement but the availability of a label-free signal for it.

\subsection{Where the Policies Are Not Expressible}

Two architectural observations bound the scope further, and neither is
an experimental outcome.

No widely used audio--language model presents a frequency axis to the
language model. Whisper-derived encoders, which cover most current
systems, fold the mel axis into the channel dimension, so the token
sequence is indexed by time alone; query-based adapters pool to an
unordered set. Frequency is to audio what colour is to images: present
in the signal, absent from the token grid. Selection over a
frequency axis is therefore not expressible, independently of whether it
would help.

Vision--language--action models fix their input resolution. OpenVLA
resizes every observation to $224\times224$ with patch-14 towers and
emits exactly $256$ visual tokens regardless of input, and SmolVLA pads
to a fixed size. Without a resolution ladder there is no budget to
reallocate, so AVR has no free parameter in these systems. The same
property makes cropping automatically budget-matched, which removes the
equal-compute confound that otherwise dominates such comparisons.

\subsection{Summary}

The measurement transfers: a depth at which a modality stops being read
exists in every modality we tested, and the depth of peak sensitivity is
nearly constant across them. The policies transfer under stated
conditions: AVR requires a knob that restores destroyed information, and
the boundary depth bounds how much budget it can reinvest through
Equation~\ref{eq:kappa}; placement requires both room above the
reference and a label-free signal that finds it. These conditions were
derived from the vision experiments and predicted each cross-modal
outcome in advance, including the negative ones.
